% Converted from README.md
\subsection{C\# - "imitatori" baza podataka}\label{c---imitatori-baza-podataka}

Objekti imitatori, iako korisni, ne mogu zameniti ponašanje pravih sistema za upravljanje bazama podataka. Iako je testove potrebno izvršiti nad produkcijskom bazom podataka pre isporučivanja aplikacije, neefikasno je te testove pokretati tokom razvoja aplikacije. Želeli bismo da testove pokrenemo u okruženju što sličnijem pravoj bazi podataka, ali ne bismo da gubimo na efikasnosti - potrebno je bazu pripremiti za svaki test posebno, što uključuje veliki broj upita. S druge strane, možda smo već napisali imitatore i zadovoljni smo ponašanjem naših servisa, ali bismo da testiramo logiku ostvarivanja odnosno raskidanja veze sa bazom. U ovakvim slučajevima, ali i mnogim drugim, izvršavanje testova nad "pravom" bazom podataka je nešto što bismo voleli da imamo. Međutim, voleli bismo da to sve bude transparentno i da ne zahteva nikakve izmene u kodu, ali i takođe dovoljno efikasno tako da ne odlazi dosta vremena na izvršavanje testova.

Objektno-relacioni maperi (engl. Object Relational Mappers, skr. ORM) često se koriste radi apstrahovanja specifičnosti konkretnog sistema za upravljanje bazama podataka. Neki poznati ORM radni okviri su \href{https://hibernate.org/}{Hibernate} (Java), \href{https://learn.microsoft.com/en-us/ef/}{Entity Framework} (.NET). U ovom primeru ćemo iskoristiti pogodnosti koje Entity Framework i SUBP \href{https://www.sqlite.org/}{SQLite} pružaju, kako bismo testove izvršili nad privremenom bazom podataka u radnoj memoriji.

Postoje sistemi za upravljanje bazama podataka koji mogu da rade sa privremenim bazama podataka skladištenim u radnoj memoriji. Jedan od takvih SUBP je \href{https://www.sqlite.org/}{SQLite}, koji pruža tzv. \emph{in-memory database provider} upravo za ove svrhe. SQLite čuva baze podataka u datotekama na sistemu datoteka. Da bismo kreirali bazu podataka u memoriji, potrebno je da u nisku za konekciju na bazu ubacimo \texttt{:memory:} na mesto gde bi inače išla putanja do datoteke na sistemu datoteka. Privremena baza živi u radnoj memoriji sve dok postoji otvorena veza ka njoj. Raskid veze zato briše bazu podataka, što odgovara testovima koji treba da se pokreću u izolovanom okruženju.

Klasa \texttt{SampleDbContext} predstavlja kontekst veze s bazom podataka. Taj kontekst u sebi ima svojstva tipa \texttt{DbSet\textless{}T\textgreater{}} koja će se mapirati u odgovarajuće tabele. Pošto koristimo Entity Framework, možemo ga konfigurisati tako da, na osnovu konfiguracije koju korisnik navodi, koristimo odgovarajući zadnji deo (engl. back-end) koji će komunicirati sa odgovarajućim SUBP. Recimo da želimo da podržimo naredne SUBP:

\begin{minted}{csharp}
public enum Provider
{
    Sqlite = 0,
    PostgreSql = 1,
    SqlServer = 2,
    SqliteInMemory = 3
}
\end{minted}

Prilikom konfigurisanja EF radnog okvira, možemo odabrati odgovarajući zadnji deo:

\begin{minted}{csharp}
switch (this.Provider) {
    case Provider.PostgreSql:
        optionsBuilder.UseNpgsql(this.ConnectionString);
        break;
    case Provider.Sqlite:
    case Provider.SqliteInMemory:
        optionsBuilder.UseSqlite(this.ConnectionString);
        break;
    case Provider.SqlServer:
        optionsBuilder.UseSqlServer(this.ConnectionString);
        break;
    default:
        throw new NotSupportedException("Provider not supported!");
}
\end{minted}

Možemo implementirati provajder baze za testove tako što prosledimo odgovarajuću nisku za konekciju ka bazi podataka:

\begin{minted}{csharp}
public class TestDbProvider
{
    public readonly string ConnectionString { get; }
    public readonly SqliteConnection DatabaseConnection { get; }

    public TestDbProvider()
    {
        ConnectionString = "DataSource=:memory:;foreign keys=true;";
        DatabaseConnection = new SqliteConnection(ConnectionString);
    }

    private void CreateContext()
    {
        var options = new DbContextOptionsBuilder<SampleDbContext>()
            .UseSqlite(DatabaseConnection)
            .Options;

        return new SampleDbContext(
            SampleDbContext.Provider.SqliteInMemory,
            ConnectionString,
            options
        );
    }

    // ...
}
\end{minted}

Za izvršavanje testova je neophodno da:

\begin{itemize}
\tightlist
\item
  ostvarimo konekciju ka bazi
\item
  ubacimo podatke u bazu (pošto se baza uvek briše nakon raskidanja konekcije)
\item
  odradimo logiku koju test treba da proveri
\item
  proverimo rezultat
\item
  raskinemo vezu sa bazom
\end{itemize}

Da bismo smanjili ponavljanje koda, dodaćemo metod \texttt{SetupAlterAndVerify} koji će da primi funkcije:

\begin{itemize}
\tightlist
\item
  \texttt{void\ Setup(SampleDbContext\ ctx)} - priprema bazu podataka za test
\item
  \texttt{void\ Alter(SampleDbContext\ ctx)} - izvršava logiku koju test treba da proveri
\item
  \texttt{void\ Verify(SampleDbContext\ ctx)} - testira rezultujuće stanje baze podataka
\end{itemize}

\begin{minted}{csharp}

public void SetupAlterAndVerify(
    Action<SampleDbContext>? setup,
    Action<SampleDbContext>? alter,
    Action<SampleDbContext>? verify)
{
    DatabaseConnection.Open();
    try {
        this.CreateDatabase();
        this.SeedDatabase();

        if (setup is not null) {
            using SampleDbContext context = this.CreateContext();
            setup(context);
            context.SaveChanges();
        }

        if (alter is not null) {
            using SampleDbContext context = this.CreateContext();
            alter(context);
            context.SaveChanges();
        }

        if (verify is not null) {
            using SampleDbContext context = this.CreateContext();
            verify(context);
        }
    } finally {
        DatabaseConnection.Close();
    }
}
\end{minted}

Primetimo da, iako kontekst kreiramo više puta, veza ka bazi i dalje ostaje aktivna dok se ne pozove \texttt{DatabaseConnection.Close()} metod. Kreiranje zasebnih konteksta je poželjno pošto bismo da sačuvamo stanje baze nakon svakog koraka (setup, alter, verify). Tip \texttt{Action\textless{}T1,\ T2,\ ...,\ Tn\textgreater{}} u programskom jeziku C\# predstavlja funkciju: \texttt{void\ f(T1,\ T2,\ ...\ Tn)}. Oznaka \texttt{?} je skraćenica za tip \texttt{Nullable\textless{}T\textgreater{}} koji predstavlja opcioni tip. Drugim rečima, metodi \texttt{SetupAlterAndVerify} prosleđujemo opcione akcije, i možemo da odlučimo da naš test ne mora da ima neku od njih (tako što prosledimo \texttt{null}, stoga provere u telu funkcije pre poziva funkcija \texttt{setup}, \texttt{alter} i \texttt{verify}). Ključna reč \texttt{using} je deo upravljanja resursa nad objektima koji implementiraju \texttt{IDisposable} interfejs u programskom jeziku C\#, sa ciljem da se automatski počisti objekat nakon što kontrola toka izađe iz opsega u kojem je vidljiv (nešto nalik na \emph{try-with-resources} šablon u programskom jeziku Java). Drugim rečima, automatski će se pozvati metod \texttt{IDisposable.Dispose()} nad kontekstom koji je definisan naredbom koja je kvalifikovana ključnom rečju \texttt{using} (interfejs \texttt{IDisposable} je implementiran u natklasi klase \texttt{SampleDbContext} koja dolazi iz EF radnog okvira).

Testove onda možemo veoma jednostavno pisati:

\codefile{SampleDbTests.cs}{csharp}

Za poznavaoce asinhronog šablona zasnovanog na zadacima (engl. \href{https://learn.microsoft.com/en-us/dotnet/standard/asynchronous-programming-patterns/task-based-asynchronous-pattern-tap}{Task-Based Asynchronous Pattern, skr. TAP}) u programskom jeziku C\#, mogu od značaja biti i asinhrona varijanta metoda \texttt{SetupAlterAndVerify}:

\begin{minted}{csharp}
public async Task SetupAlterAndVerifyAsync(
    Func<SampleDbContext, Task>? setup,
    Func<SampleDbContext, Task>? alter,
    Func<SampleDbContext, Task>? verify)
{
    DatabaseConnection.Open();
    try {
        this.CreateDatabase();
        this.SeedDatabase();

        if (setup is not null) {
            await using SampleDbContext context = this.CreateContext();
            await setup(context);
            await context.SaveChangesAsync();
        }

        if (alter is not null) {
            await using SampleDbContext context = this.CreateContext();
            await alter(context);
            await context.SaveChangesAsync();
        }

        if (verify is not null) {
            await using SampleDbContext context = this.CreateContext();
            await verify(context);
        }
    } finally {
        DatabaseConnection.Close();
    }
}
\end{minted}
